\makeatletter
\def\input@path{{./}{../}{../../}{../../../}{../../../../}{preamble/}{../preamble/}{../../preamble/}{../../../preamble/}}
\makeatother

\input{preamble_exams.tex}



\input{graphs/G11_C6C7_graphs.tex}
\usepackage{longtable}

\setlength{\DistributionGraphWidth}{0.45\textwidth}
\setlength{\DistributionGraphHeight}{4cm}

\title{Grade 11 -- Term 1 \\ C6--C7 Practice Activity with Solutions}
\author{}
\date{}

\pagestyle{empty}

% ============================================================
% GRAPH WRAPPER
% ============================================================

\begin{document}

	\maketitle
	\setcounter{tocdepth}{1}
	\setcounter{secnumdepth}{2}
	\phantomsection\label{toc}
	\tableofcontents
	
	\section*{Evaluated criterion}
	\begin{itemize}
		\item [\textbf{C6.}] Estimates interval probabilities using relative frequencies from generated samples and compares them with theoretical probabilities.
		
		\item [\textbf{C7.}] Calculates sample means and variances for different sample sizes and compares them with the theoretical mean and variance.
		
		
	\end{itemize}
	
	\vspace{-4mm}
	
% ==========================================================================
% C6 INTRODUCTION
% ==========================================================================

\section{Introduction}\label{sec:c6-introduction}

\SectionProblemsTableOfContents{
	\SectionProblemLink{sec:c6-intro-samples-histograms}{Generated samples, sample sizes, and histograms} \\
	\SectionProblemLink{sec:c6-intro-distributions}{Continuous distributions, theoretical and empirical probabilities} \\
	\SectionProblemLink{sec:c6-intro-sample-mean-variance}{Sample mean and variance from generated data}
}
\vspace{8mm}

Criterion C6 focuses on estimating interval probabilities from generated
samples and comparing these empirical estimates with the theoretical
probabilities predicted by a continuous probability distribution.

The comparison can be understood through two representations of the same
distribution. The probability density function (PDF) represents the
theoretical model as a smooth curve. A histogram represents the observations
obtained from a finite generated sample. The theoretical probability of an
interval is determined by the area under the PDF, while the empirical
probability is estimated from the relative frequency of observations in the
corresponding interval.

The normal distribution provides a useful example because its theoretical
density has a smooth bell-shaped form, while a histogram generated from the
same distribution produces a finite and irregular approximation to that
shape.


% ==========================================================================
% CONTINUOUS PROBABILITY DISTRIBUTIONS
% ==========================================================================
\subsection{Real-world data, probability distributions, and histograms}
\label{sec:c6-intro-samples-histograms}

Real-world data often show patterns of variation that can be described using
probability distributions. A theoretical distribution provides a model for
how values of a variable may be distributed, while a dataset contains the
observations actually collected or generated from that process.

A dataset with $n$ observations can be written as

$$
x_1,x_2,\ldots,x_n.
$$

The observations can be organized into a histogram, which groups values into
intervals and shows their frequencies. The histogram provides an empirical
description of the data that can be compared with the theoretical probability
distribution.

\begin{center}
	\begin{tabular}{ccc}
		\begin{minipage}[t]{0.30\textwidth}
			\centering
			\CSampleCountHistogramTwenty
			\captionof{figure}{
				Count histogram from a sample of $n=20$ observations.
			}
			\label{fig:c6-sample-count-histogram-20}
		\end{minipage}
		&
		\begin{minipage}[t]{0.30\textwidth}
			\centering
			\CSampleCountHistogramHundred
			\captionof{figure}{
				Count histogram from a sample of $n=100$ observations.
			}
			\label{fig:c6-sample-count-histogram-100}
		\end{minipage}
		&
		\begin{minipage}[t]{0.30\textwidth}
			\centering
			\CSampleCountHistogramThousand
			\captionof{figure}{
				Count histogram from a sample of $n=1000$ observations.
			}
			\label{fig:c6-sample-count-histogram-1000}
		\end{minipage}
	\end{tabular}
\end{center}

\begin{center}
	\begin{tabular}{ccc}
		\begin{minipage}[t]{0.30\textwidth}
			\centering
			\CSampleHistogramTwenty
			\captionof{figure}{
				Histogram from a sample of $n=20$ observations.
			}
			\label{fig:c6-sample-histogram-20}
		\end{minipage}
		&
		\begin{minipage}[t]{0.30\textwidth}
			\centering
			\CSampleHistogramHundred
			\captionof{figure}{
				Histogram from a sample of $n=100$ observations.
			}
			\label{fig:c6-sample-histogram-100}
		\end{minipage}
		&
		\begin{minipage}[t]{0.30\textwidth}
			\centering
			\CSampleHistogramThousand
			\captionof{figure}{
				Histogram from a sample of $n=1000$ observations.
			}
			\label{fig:c6-sample-histogram-1000}
		\end{minipage}
	\end{tabular}
\end{center}

Different samples from the same process can produce different histograms.
Larger samples generally provide a more stable representation of the
underlying distribution because random variation has less influence on the
overall shape.

The theoretical distribution and the observed data can also be compared
numerically. The theoretical distribution has mean $\mu$ and variance
$\sigma^2$, while the sample has mean $\bar{x}$ and variance $s^2$:

$$
\boxed{
	\begin{array}{c|c}
		\textbf{Theoretical model} & \textbf{Observed data} \\ \hline
		\text{Probability distribution}
		&
		\text{Observations }x_1,x_2,\ldots,x_n
		\\[1mm]
		\text{Theoretical mean }\mu
		&
		\text{Sample mean }\bar{x}
		\\[1mm]
		\text{Theoretical variance }\sigma^2
		&
		\text{Sample variance }s^2
\end{array}}
$$

As the sample size increases, the sample statistics generally become more
stable and tend to approximate the corresponding theoretical parameters:

$$
\bar{x}\approx\mu,
\qquad
s^2\approx\sigma^2.
$$

Thus, the relationship can be summarized as

$$
\boxed{
	\text{Real-world process}
	\;\longrightarrow\;
	\text{Observed data}
	\;\longrightarrow\;
	\text{Histogram and sample statistics}
	\;\longleftrightarrow\;
	\text{Probability distribution}
}
$$

The following subsection presents the equations used to calculate the sample
mean and variance from the data.

\subsection{Continuous distributions, theoretical and empirical probabilities}
\label{sec:c6-intro-distributions}

\begin{center}
	\CNormalPDFAndHistogram{0}{1}
	\captionof{figure}{
		Theoretical normal probability density function and a histogram
		from a sample generated from the same distribution. The smooth
		curve represents the theoretical model, while the bars represent
		the observed sample.
	}
	\label{fig:c6-normal-pdf-histogram}
\end{center}

Continuous probability distributions provide theoretical models for
continuous random variables. For an interval $[a,b]$, the theoretical
probability is determined by the area under the probability density function:

$$
P(a\leq X\leq b)
=
\int_a^b f(x)\,dx.
$$

When a finite sample is generated from the distribution, the corresponding
probability can instead be estimated from the observations. If $k$ of the
$n$ observations fall within the interval, the empirical probability is
given by the relative frequency

$$
\widehat{P}(a\leq X\leq b)
=
\frac{k}{n},
\qquad
k=\#\{x_i:a\leq x_i\leq b\}.
$$

The distinction between the two probabilities can be summarized as follows:

$$
\boxed{
	\begin{array}{c|c}
		\textbf{Theoretical probability} & \textbf{Empirical probability} \\ \hline
		\text{Based on the probability distribution}
		&
		\text{Based on a generated sample}
		\\[1mm]
		\text{Area under the PDF}
		&
		\text{Relative frequency}
		\\[1mm]
		\displaystyle
		P(a\leq X\leq b)
		=
		\int_a^b f(x)\,dx
		&
		\displaystyle
		\widehat{P}(a\leq X\leq b)
		=
		\frac{k}{n}
\end{array}}
$$

The following pair of graphs illustrates the two calculations separately:
\vspace{-4mm}
\begin{center}
	\begin{tabular}{cc}
		\begin{minipage}[t]{0.48\textwidth}
			\centering
			\CNormalPDFGraph{0}{1}
			\captionof{figure}{
				Theoretical probability density function of a normal
				distribution. The area under the curve over an interval
				represents the theoretical probability of that interval.
			}
			\label{fig:c6-normal-pdf}
		\end{minipage}
		&
		\begin{minipage}[t]{0.48\textwidth}
			\centering
			\CSampleHistogramNormal{0}{1}
			\captionof{figure}{
				Histogram of a finite sample generated from the same
				distribution. The bars show the observed frequencies
				within the histogram bins.
			}
			\label{fig:c6-normal-histogram}
		\end{minipage}
	\end{tabular}
\end{center}

The histogram is useful because it connects the generated observations with
the empirical calculation. The bars show where the observations occur, and
the observations within the selected interval determine $k$.

For all the continuous distributions used in this criterion, the method for
calculating the empirical probability is the same:

$$
\widehat{P}
=
\frac{\text{number of observations in the interval}}
{\text{total number of observations}}.
$$

The method for calculating the theoretical probability depends on the
probability density function of the distribution.

The key comparison is therefore between the area represented by the
theoretical PDF and the proportion of observations represented by the
histogram:

$$
\boxed{
	\text{PDF area}
	\quad\longleftrightarrow\quad
	\text{histogram relative frequency}
}
$$

For a sufficiently large generated sample, these two quantities should
generally be close, while differences are expected because the histogram is
based on a finite random sample.

The difference between the empirical and theoretical probabilities is

$$
\text{Difference}
=
\widehat{P}(a\leq X\leq b)
-
P(a\leq X\leq b).
$$

The absolute difference is

$$
\text{Absolute difference}
=
\left|
\widehat{P}(a\leq X\leq b)
-
P(a\leq X\leq b)
\right|.
$$

When the theoretical probability is nonzero, the percentage error is

$$
\text{Percentage error}
=
\frac{
	\left|
	\widehat{P}(a\leq X\leq b)
	-
	P(a\leq X\leq b)
	\right|
}{
	P(a\leq X\leq b)
}
\times100\%.
$$

These quantities provide a common reference for comparing theoretical
probabilities obtained from a continuous distribution with empirical
probabilities obtained from generated samples.

\secInicio[sec:c6-introduction][sec:c6-intro-distributions]

\subsection{Sample mean and variance from generated data}
\label{sec:c6-intro-sample-mean-variance}

A theoretical probability distribution describes a model, while a finite
generated sample provides observations from that model. The theoretical
mean and variance are denoted by $\mu$ and $\sigma^2$, while the corresponding
sample statistics are $\bar{x}$ and $s^2$.

$$
\boxed{
	\begin{array}{c|c}
		\textbf{Theoretical model} & \textbf{Generated sample} \\ \hline
		\text{Probability distribution} & \text{Finite dataset}\\
		\text{Mean } \mu & \text{Sample mean } \bar{x}\\
		\text{Variance } \sigma^2 & \text{Sample variance }s^2\\
		\text{Standard deviation } \sigma & \text{Sample standard deviation }s
\end{array}}
$$
For ungrouped observations $x_1,x_2,\ldots,x_n$, the sample mean, variance,
and standard deviation are

$$
\bar{x}
=
\frac{1}{n}\sum_{i=1}^{n}x_i,
\qquad
s^2
=
\frac{1}{n-1}\sum_{i=1}^{n}(x_i-\bar{x})^2,
\qquad
s=\sqrt{s^2}.
$$

The mean describes the center of the observed data, while the variance
measures its variability around $\bar{x}$.

When observations are grouped into intervals, let $f_i$ be the frequency of
class $i$ and $m_i$ its midpoint:

$$
m_i=\frac{a_i+b_i}{2},
\qquad
n=\sum_i f_i.
$$

The grouped-data mean and variance are

$$
\bar{x}
=
\frac{\sum_i f_i m_i}{n},
\qquad
s^2
=
\frac{\sum_i f_i(m_i-\bar{x})^2}{n-1}.
$$

Grouped calculations are approximate because all observations within an
interval are represented by its midpoint.

$$
\boxed{
	\begin{array}{c|c}
		\textbf{Ungrouped data} & \textbf{Grouped data}\\ \hline
		x_1,x_2,\ldots,x_n & \text{Intervals and frequencies}\\
		\text{Exact observations} & \text{Class midpoints }m_i\\
		\displaystyle
		\bar{x}=\frac{1}{n}\sum_i x_i
		&
		\displaystyle
		\bar{x}=\frac{\sum_i f_i m_i}{n}
		\\[3mm]
		\displaystyle
		s^2=\frac{1}{n-1}\sum_i(x_i-\bar{x})^2
		&
		\displaystyle
		s^2=\frac{\sum_i f_i(m_i-\bar{x})^2}{n-1}
\end{array}}
$$

For a finite random sample generated from a theoretical distribution,

$$
\bar{x}\approx\mu,
\qquad
s^2\approx\sigma^2.
$$

These are approximations: different samples can produce different values.
As the sample size increases, the sample statistics generally become more
stable and tend toward the corresponding theoretical values.

% ==========================================================================
% C6-ONLY PRACTICE
% ==========================================================================

\section{C6 Practice: Empirical and Theoretical Probabilities}
\label{sec:c6-only-practice}

\SectionProblemsTableOfContents{
	\SectionProblemLink{sec:c6-only-raw}{Monthly household electricity use (raw data)} \\
	\SectionProblemLink{sec:c6-only-frequency}{Waiting times (frequency data)} \\
	\SectionProblemLink{sec:c6-only-grouped}{Package masses (grouped data)}
}\vspace{5mm}

The following problems focus only on estimating interval probabilities from
samples. The larger samples make the relative-frequency estimates more stable,
while the raw, frequency, and grouped formats provide different ways of
organizing the observations.

\subsection{Monthly household electricity use (raw data)}
\label{sec:c6-only-raw}

A utility company generates the following sample of monthly household
electricity use (in hundreds of kWh):

\[
\begin{array}{cccccccccc}
	3.8 & 4.1 & 4.5 & 5.2 & 5.7 & 6.0 & 6.3 & 6.8 & 7.1 & 7.5 \\
	4.0 & 4.4 & 4.9 & 5.4 & 5.9 & 6.1 & 6.5 & 6.9 & 7.3 & 8.0
\end{array}
\]

For the theoretical model,
$P(5\leq X\leq7)=0.62$.

\textbf{Question.}
\begin{enumerate}
	\CriterionTask{subsec:c6-only-raw-solution}{6}{
		Estimate $P(5\leq X\leq7)$ from the sample. Then calculate the
		absolute error and the percentage error relative to the theoretical
		probability.
	}
\end{enumerate}

\subsubsection*{C6 Solution}
\phantomsection\label{subsec:c6-only-raw-solution}

There are $n=20$ observations. The values in the interval from $5$ to $7$,
inclusive, are

\[
5.2,\ 5.7,\ 6.0,\ 6.3,\ 6.8,\ 5.4,\ 5.9,\ 6.1,\ 6.5,\ 6.9,
\]

so $k=10$. Therefore,

\[
\widehat P(5\leq X\leq7)=\frac{k}{n}=\frac{10}{20}=0.50.
\]

The absolute error is

\[
E=|0.50-0.62|=0.12,
\]

and the percentage error is

\[
\frac{|0.50-0.62|}{0.62}\times100\%
=\frac{0.12}{0.62}\times100\%
\approx19.35\%.
\]

Thus the empirical probability is $0.50$, which is $0.12$ below the
theoretical probability.

\subsection{Waiting times (frequency data)}
\label{sec:c6-only-frequency}

A clinic records the waiting times of $120$ patients in the following
frequency table:

\[
\begin{array}{c|cccccc}
	\text{Waiting time (min)} & 5 & 10 & 15 & 20 & 25 & 30 \\ \hline
	\text{Frequency} & 8 & 20 & 34 & 30 & 18 & 10
\end{array}
\]

For the theoretical model,
$P(10\leq X\leq20)=0.72$.

\textbf{Question.}
\begin{enumerate}
	\CriterionTask{subsec:c6-only-frequency-solution}{6}{
		Estimate the probability that a patient waits from $10$ to $20$
		minutes, inclusive. Calculate its absolute error and compare the
		empirical estimate with the theoretical probability.
	}
\end{enumerate}

\subsubsection*{C6 Solution}
\phantomsection\label{subsec:c6-only-frequency-solution}

The total frequency is

\[
n=8+20+34+30+18+10=120.
\]

The waiting times $10$, $15$, and $20$ minutes are in the interval, so

\[
k=20+34+30=84.
\]

Hence

\[
\widehat P(10\leq X\leq20)=\frac{84}{120}=0.70.
\]

The absolute error is

\[
E=|0.70-0.72|=0.02.
\]

Thus the empirical probability is $0.70$, which is $0.02$ below the
theoretical probability of $0.72$.

\subsection{Package masses (grouped data)}
\label{sec:c6-only-grouped}

A distribution center groups the masses of $200$ packages as follows:

\[
\begin{array}{c|ccccc}
	\text{Package mass (kg)}
	& 0\text{--}2 & 2\text{--}4 & 4\text{--}6 & 6\text{--}8 & 8\text{--}10 \\ \hline
	\text{Frequency} & 18 & 46 & 72 & 44 & 20
\end{array}
\]

The classes are left-inclusive and right-exclusive. For the theoretical
model, $P(2\leq X<8)=0.83$.

\textbf{Question.}
\begin{enumerate}
	\CriterionTask{subsec:c6-only-grouped-solution}{6}{
		Estimate $P(2\leq X<8)$ from the grouped sample. Calculate the
		absolute error and the percentage error relative to the theoretical
		probability.
	}
\end{enumerate}

\subsubsection*{C6 Solution}
\phantomsection\label{subsec:c6-only-grouped-solution}

The total frequency is

\[
n=18+46+72+44+20=200.
\]

The interval $2\leq X<8$ contains the three classes $2$--$4$, $4$--$6$,
and $6$--$8$. Their combined frequency is

\[
k=46+72+44=162.
\]

Therefore,

\[
\widehat P(2\leq X<8)=\frac{162}{200}=0.81.
\]

The absolute error is

\[
E=|0.81-0.83|=0.02,
\]

and the percentage error is

\[
\frac{|0.81-0.83|}{0.83}\times100\%
=\frac{0.02}{0.83}\times100\%
\approx2.41\%.
\]

Thus the empirical estimate is $0.81$, which is close to, but slightly
below, the theoretical probability of $0.83$.

% ==========================================================================
% C7-ONLY PRACTICE
% ==========================================================================

\section{C7 Practice: Sample Mean, Variance, and Standard Deviation}
\label{sec:c7-only-practice}

\SectionProblemsTableOfContents{
	\SectionProblemLink{sec:c7-only-raw}{Machine calibration errors (raw data)} \\
	\SectionProblemLink{sec:c7-only-frequency}{Books borrowed (frequency data)} \\
	\SectionProblemLink{sec:c7-only-grouped}{Commute times (grouped data)}
}\vspace{5mm}

The samples in this section are deliberately small so that the sample mean,
sample variance, and sample standard deviation can be calculated clearly by
hand. Recall that the sample-variance denominator is $n-1$.

\subsection{Machine calibration errors (raw data)}
\label{sec:c7-only-raw}

The calibration errors (in tenths of a millimeter) from five trials are

\[
4,\quad5,\quad6,\quad7,\quad8.
\]

The theoretical model has $\mu=6$ and $\sigma^2=2.25$.

\textbf{Question.}
\begin{enumerate}
	\CriterionTask{subsec:c7-only-raw-solution}{7}{
		Calculate the sample mean, sample variance, and sample standard
		deviation. Compare the sample mean and variance with the theoretical
		values.
	}
\end{enumerate}

\subsubsection*{C7 Solution}
\phantomsection\label{subsec:c7-only-raw-solution}

There are $n=5$ observations. First, calculate the sample mean:

\[
\bar{x}
=\frac{4+5+6+7+8}{5}
=\frac{30}{5}
=6.
\]

Next, calculate the squared deviations from $\bar{x}=6$:

\[
\begin{array}{c|ccccc}
	x_i & 4 & 5 & 6 & 7 & 8 \\ \hline
	x_i-\bar{x} & -2 & -1 & 0 & 1 & 2 \\
	(x_i-\bar{x})^2 & 4 & 1 & 0 & 1 & 4
\end{array}
\]

Their sum is $4+1+0+1+4=10$. Therefore, the sample variance is

\[
s^2
=\frac{\sum_{i=1}^{5}(x_i-\bar{x})^2}{n-1}
=\frac{10}{5-1}
=2.50.
\]

The sample standard deviation is

\[
s=\sqrt{s^2}=\sqrt{2.50}\approx1.58.
\]

Thus,

\[
\boxed{\bar{x}=6=\mu,\qquad
s^2=2.50\approx\sigma^2=2.25,\qquad s\approx1.58.}
\]

\subsection{Books borrowed (frequency data)}
\label{sec:c7-only-frequency}

The numbers of books borrowed by four students are summarized below:

\[
\begin{array}{c|ccc}
	x_i & 2 & 4 & 6 \\ \hline
	f_i & 1 & 2 & 1
\end{array}
\]

The theoretical model has $\mu=4$ and $\sigma^2=3$.

\textbf{Question.}
\begin{enumerate}
	\CriterionTask{subsec:c7-only-frequency-solution}{7}{
		Calculate the sample mean, sample variance, and sample standard
		deviation from the frequency table. Compare the sample mean and
		variance with the theoretical values.
	}
\end{enumerate}

\subsubsection*{C7 Solution}
\phantomsection\label{subsec:c7-only-frequency-solution}

The total frequency, and hence the sample size, is

\[
n=1+2+1=4.
\]

The sample mean is

\[
\bar{x}
=\frac{\sum_i f_i x_i}{n}
=\frac{1(2)+2(4)+1(6)}{4}
=\frac{16}{4}
=4.
\]

Using $\bar{x}=4$, calculate each weighted squared deviation:

\[
\begin{array}{c|ccc}
	x_i & 2 & 4 & 6 \\ \hline
	f_i & 1 & 2 & 1 \\
	(x_i-\bar{x})^2 & 4 & 0 & 4 \\
	f_i(x_i-\bar{x})^2 & 4 & 0 & 4
\end{array}
\]

The sum of the weighted squared deviations is $4+0+4=8$. Thus,

\[
s^2
=\frac{\sum_i f_i(x_i-\bar{x})^2}{n-1}
=\frac{8}{4-1}
=\frac83
\approx2.67.
\]

The sample standard deviation is

\[
s=\sqrt{s^2}=\sqrt{\frac83}\approx1.63.
\]

Therefore,

\[
\boxed{\bar{x}=4=\mu,\qquad
s^2\approx2.67\approx\sigma^2=3,\qquad s\approx1.63.}
\]

\subsection{Commute times (grouped data)}
\label{sec:c7-only-grouped}

The commute times of five employees are grouped into three classes:

\[
\begin{array}{c|ccc}
	\text{Commute time (min)} & 0\text{--}10 & 10\text{--}20 & 20\text{--}30 \\ \hline
	\text{Frequency} & 1 & 3 & 1
\end{array}
\]

The theoretical model has $\mu=15$ minutes and
$\sigma^2=48$ minutes$^2$.

\textbf{Question.}
\begin{enumerate}
	\CriterionTask{subsec:c7-only-grouped-solution}{7}{
		Use class midpoints to estimate the sample mean, sample variance, and
		sample standard deviation. Compare the estimated mean and variance
		with the theoretical values.
	}
\end{enumerate}

\subsubsection*{C7 Solution}
\phantomsection\label{subsec:c7-only-grouped-solution}

First, calculate each class midpoint:

\[
m_1=\frac{0+10}{2}=5,\qquad
m_2=\frac{10+20}{2}=15,\qquad
m_3=\frac{20+30}{2}=25.
\]

The total frequency is

\[
n=1+3+1=5.
\]

Using the midpoints as representative values, the estimated sample mean is

\[
\bar{x}
=\frac{\sum_i f_i m_i}{n}
=\frac{1(5)+3(15)+1(25)}{5}
=\frac{75}{5}
=15\text{ min}.
\]

Next, calculate the weighted squared deviations:

\[
\begin{array}{c|ccc}
	m_i & 5 & 15 & 25 \\ \hline
	f_i & 1 & 3 & 1 \\
	(m_i-\bar{x})^2 & 100 & 0 & 100 \\
	f_i(m_i-\bar{x})^2 & 100 & 0 & 100
\end{array}
\]

Their sum is $100+0+100=200$. The estimated sample variance is

\[
s^2
=\frac{\sum_i f_i(m_i-\bar{x})^2}{n-1}
=\frac{200}{5-1}
=50\text{ min}^2.
\]

The estimated sample standard deviation is

\[
s=\sqrt{s^2}=\sqrt{50}\approx7.07\text{ min}.
\]

Therefore,

\[
\boxed{\bar{x}=15\text{ min}=\mu,\qquad
s^2=50\text{ min}^2\approx\sigma^2=48\text{ min}^2,\qquad
s\approx7.07\text{ min}.}
\]

These results are estimates because all observations in a class are
represented by the class midpoint.

\section{Samples, Empirical Probabilities, and Sample Statistics}
\label{sec:c6-specific-statistics}

\SectionProblemsTableOfContents{
	\SectionProblemLink{sec:c6c7-data-raw}{Investment returns (raw data)} \\
	\SectionProblemLink{sec:c6-probability-frequency}{Daily energy consumption (frequency data)} \\
	\SectionProblemLink{sec:c6-probability-grouped}{Delivery times (grouped data)}
}\vspace{5mm}

The following examples use small samples so that each step of the calculations can be shown clearly. In real-world applications, sample sizes are typically much larger to obtain more reliable estimates.

\subsection{Investment returns (raw data)}
\label{sec:c6c7-data-raw}

A financial analyst records these returns (in percent) over six days:

\[
2,\quad 2,\quad 3,\quad 3,\quad 4,\quad 4.
\]

The theoretical model has $\mu=3$, $\sigma^2=1$, and
$P(2\leq X\leq4)=0.80$.

\textbf{Question.}
\begin{enumerate}
	\CriterionTask{subsec:c6-probability-raw-b}{6}{
		Calculate the empirical probability of a return between $2\%$ and
		$4\%$, inclusive, and its absolute error.
	}
	\CriterionTask{subsec:c7-meanvar-raw-b}{7}{
		Calculate the sample mean and sample variance, and compare them with
		the theoretical values.
	}
\end{enumerate}

\subsubsection*{C6 Solution}
\phantomsection\label{subsec:c6-probability-raw-b}

All six observations are in the interval, so $n=6$ and $k=6$. Hence

\[
\widehat P(2\leq X\leq4)=\frac{k}{n}=\frac{6}{6}=1.
\]

The absolute error is

\[
E=|\widehat P-P|=|1-0.80|=0.20.
\]

Thus the empirical probability is $1$, which is $0.20$ above the
theoretical probability.

\subsubsection*{C7 Solution}
\phantomsection\label{subsec:c7-meanvar-raw-b}

The sample mean is

\[
\bar{x}
=\frac{2+2+3+3+4+4}{6}
=\frac{18}{6}
=3.
\]

Using $\bar{x}=3$, the sample variance is

\[
\begin{aligned}
s^2
&=\frac{(2-3)^2+(2-3)^2+(3-3)^2+(3-3)^2+(4-3)^2+(4-3)^2}{6-1}\\
&=\frac{1+1+0+0+1+1}{5}
=\frac45
=0.80.
\end{aligned}
\]

Therefore,

\[
\boxed{\bar{x}=3=\mu,\qquad s^2=0.80\approx\sigma^2=1.}
\]

\subsection{Daily energy consumption (frequency data)}
\label{sec:c6-probability-frequency}

A facility records daily consumption (in MWh) as follows:

\[
\begin{array}{c|ccc}
	x_i & 8 & 10 & 12 \\ \hline
	f_i & 1 & 2 & 1
\end{array}
\]

The theoretical model has $\mu=10$, $\sigma^2=3$, and
$P(10\leq X\leq12)=0.60$.

\textbf{Question.}
\begin{enumerate}
	\CriterionTask{subsec:c6-probability-frequency-b}{6}{
		Calculate the empirical probability of consumption between $10$ and
		$12$ MWh, inclusive, and its absolute error.
	}
	\CriterionTask{subsec:c7-meanvar-frequency-b}{7}{
		Calculate the sample mean and sample variance, and compare them with
		the theoretical values.
	}
\end{enumerate}

\subsubsection*{C6 Solution}
\phantomsection\label{subsec:c6-probability-frequency-b}

The total frequency and the frequency in the interval are

\[
n=1+2+1=4,
\qquad
k=2+1=3.
\]

Therefore,

\[
\widehat P(10\leq X\leq12)=\frac34=0.75,
\]

and the absolute error is

\[
E=|0.75-0.60|=0.15.
\]

Thus the empirical probability is $0.75$, which is $0.15$ above the
theoretical probability.

\subsubsection*{C7 Solution}
\phantomsection\label{subsec:c7-meanvar-frequency-b}

The sample mean is

\[
\bar{x}
=\frac{1(8)+2(10)+1(12)}{4}
=\frac{40}{4}
=10.
\]

The sample variance is

\[
\begin{aligned}
s^2
&=\frac{1(8-10)^2+2(10-10)^2+1(12-10)^2}{4-1}\\
&=\frac{4+0+4}{3}
=\frac83
\approx2.67.
\end{aligned}
\]

Therefore,

\[
\boxed{\bar{x}=10=\mu,\qquad s^2\approx2.67\approx\sigma^2=3.}
\]

\subsection{Delivery times (grouped data)}
\label{sec:c6-probability-grouped}

A company groups delivery times into three classes:

\[
\begin{array}{c|ccc}
	\text{Delivery time (min)} & 10\text{--}20 & 20\text{--}30 & 30\text{--}40 \\ \hline
	\text{Frequency} & 1 & 2 & 1
\end{array}
\]

The theoretical model has $\mu=25$ minutes, $\sigma^2=75$ minutes$^2$,
and $P(20\leq X<30)=0.45$.

\textbf{Question.}
\begin{enumerate}
	\CriterionTask{subsec:c6-probability-grouped-b}{6}{
		Estimate the empirical probability of a delivery time from $20$ up
		to (but not including) $30$ minutes, and calculate its absolute error.
	}
	\CriterionTask{subsec:c7-meanvar-grouped-b}{7}{
		Estimate the sample mean and sample variance using class midpoints,
		and compare them with the theoretical values.
	}
\end{enumerate}

\subsubsection*{C6 Solution}
\phantomsection\label{subsec:c6-probability-grouped-b}

The total frequency is $n=1+2+1=4$. The selected class has frequency
$k=2$, so

\[
\widehat P(20\leq X<30)=\frac24=0.50.
\]

The absolute error is

\[
E=|0.50-0.45|=0.05.
\]

Thus the empirical probability is $0.50$, which is $0.05$ above the
theoretical probability.

\subsubsection*{C7 Solution}
\phantomsection\label{subsec:c7-meanvar-grouped-b}

The three class midpoints are

\[
15,\qquad25,\qquad35.
\]

Using these midpoints, the estimated sample mean is

\[
\bar{x}
=\frac{1(15)+2(25)+1(35)}{4}
=\frac{100}{4}
=25\text{ min}.
\]

The estimated sample variance is

\[
\begin{aligned}
s^2
&=\frac{1(15-25)^2+2(25-25)^2+1(35-25)^2}{4-1}\\
&=\frac{100+0+100}{3}
=\frac{200}{3}
\approx66.67\text{ min}^2.
\end{aligned}
\]

Therefore,

\[
\boxed{\bar{x}=25\text{ min}=\mu,\qquad
s^2\approx66.67\text{ min}^2\approx\sigma^2=75\text{ min}^2.}
\]

These grouped-data results are estimates because every observation in a
class is represented by its midpoint.


\end{document}
